\documentclass[11pt,reqno]{amsart}
\usepackage[margin=1.15in]{geometry}
\usepackage{amsmath,amssymb,amsthm}
\usepackage{hyperref}

\theoremstyle{plain}
\newtheorem{theorem}{Theorem}[section]
\newtheorem{lemma}[theorem]{Lemma}
\newtheorem{proposition}[theorem]{Proposition}
\newtheorem{corollary}[theorem]{Corollary}
\theoremstyle{definition}
\newtheorem{definition}[theorem]{Definition}
\theoremstyle{remark}
\newtheorem{remark}[theorem]{Remark}

\newcommand{\Z}{\mathbb{Z}}
\newcommand{\Q}{\mathbb{Q}}
\newcommand{\R}{\mathbb{R}}
\newcommand{\F}{\mathbb{F}}
\newcommand{\OK}{\mathcal{O}_K}
\newcommand{\Ht}{\operatorname{Ht}}
\newcommand{\vp}{v_{\mathfrak{p}}}
\newcommand{\ord}{\operatorname{ord}}
\newcommand{\p}{\mathfrak{p}}
\newcommand{\q}{\mathfrak{q}}
\newcommand{\N}{\mathrm{N}}

\begin{document}

\title[Densities of Eisenstein--Dumas polynomials and Heilbronn's criterion]
  {Densities of Eisenstein--Dumas polynomials \\ and Heilbronn's criterion}

\author{Michail Karatarakis}

\date{September 2026}

\begin{abstract}
Heilbronn's criterion is a sufficient condition for a number field with a totally ramified
prime to fail to be norm-Euclidean. Hibbler, McGown and Trevi\~no recently showed that, among
the monic integer polynomials of degree $n$ that are Eisenstein at a prime $p$, the proportion
of those to which the criterion applies has lower density at least $\max\{2/27,\,1-\varepsilon(p)\}$.
We extend their results in four directions. First, we replace the Eisenstein condition by the
much weaker Eisenstein--Dumas condition, that the Newton polygon at $p$ consist of a single
segment of slope $-m/n$ with $\gcd(m,n)=1$; we prove that this condition already forces total
ramification, and we compute the density of the resulting family to be
$(1-1/p)\,p^{-S(m,n)}$ with $S(m,n)=\tfrac12(m+1)(n+1)-1$, which for $m=1$ recovers the classical
Eisenstein density. Second, we prove a single master density theorem, indexed by a finite set $Q$
of auxiliary primes, of which the constant $2/27$ and the bound $1-\varepsilon(p)$ are the two
extreme cases; intermediate choices of $Q$ give strictly better bounds in the intermediate range
of $p$, for instance $136/675 \approx 0.2015$ for all $p \ge 225$. Third, we give a proof of
Heilbronn's criterion that uses no embeddings and no Galois closure, resting instead on the fact
that the norm form of a totally ramified prime reduces to the $n$-th power map on the residue
field. Along the way we isolate the exact non-norm statement that the argument requires, which is
slightly stronger than the one usually quoted. Finally, we make the removal of the hypothesis
$\gcd(p-1,n)=1$ effective: imposing the two congruence conditions on the auxiliary integer as a
single arithmetic progression leaves only character sums modulo the prime $p$ itself, so that the
Pólya--Vinogradov inequality applies without any reduction of imprimitive characters, and the
resulting criterion is an explicit inequality with no unspecified constant.
\end{abstract}

\maketitle

\section{Introduction}

Let $f(x)=a_nx^n+\dots+a_1x+a_0 \in \Z[x]$ and let
$\Ht(f)=\max\{|a_0|,\dots,|a_n|\}$ be the standard height, so that for every $X>0$ the set of
$f$ in a given family with $\Ht(f)\le X$ is finite. For a monic $f$ of degree $n$ and $X\ge1$ the
condition $\Ht(f)\le X$ says exactly that $|a_i|\le X$ for $0\le i\le n-1$, so the monic
polynomials of degree $n$ and height at most $X$ are counted by $(2X+1)^n$. Given a family $\mathcal F$ of polynomials and
a distinguished subfamily $\mathcal F^*$, we write $N(X)$ and $N^*(X)$ for the corresponding
counting functions and $\delta(X)=N^*(X)/N(X)$ for the proportion; the lower density of
$\mathcal F^*$ in $\mathcal F$ is $\liminf_{X\to\infty}\delta(X)$.

A number field $K$ is \emph{norm-Euclidean} if for all $\alpha,\beta\in\OK$ with $\beta\ne 0$
there exist $\gamma,\rho\in\OK$ with $\alpha=\gamma\beta+\rho$ and $|\N(\rho)|<|\N(\beta)|$.
Deciding which fields are norm-Euclidean is hard, and has been completed only in degree two; see
the survey \cite{Lemmermeyer}. When $K$ possesses a totally ramified prime there is a classical
obstruction available, due to Heilbronn \cite{H1,H2} (see also \cite{EK,Egami}).

\begin{theorem}[Heilbronn's criterion]\label{thm:heilbronn-intro}
Let $K$ be a number field of degree $n$ and let $p$ be a prime that is totally ramified in $K$.
Suppose $p=a+b$ with $a,b>0$, that $a$ is an $n$-th power residue modulo $p$, and that neither
$a$ nor $-b$ is the norm of an algebraic integer of $K$. Then $K$ is not norm-Euclidean.
\end{theorem}

Hibbler, McGown and Trevi\~no \cite{HMT} asked how often this criterion applies inside a natural
family of polynomials, and answered the question for Eisenstein polynomials. Recall that a monic
$f\in\Z[x]$ of degree $n$ is $p$-Eisenstein if $p \mid a_i$ for $i<n$ and $p^2 \nmid a_0$; such an
$f$ is irreducible, and $p$ is totally ramified in $K_f=\Q[x]/(f)$. Their main result is that for
odd $n\ge 3$ and a prime $p\ge5$ with $\gcd(p-1,n)=1$, the lower density of the $p$-Eisenstein
polynomials of degree $n$ whose field is not norm-Euclidean is at least
$\max\{2/27,\ 1-\varepsilon(p)\}$, with $\varepsilon(p)\to0$ effectively.

The Eisenstein condition is only the simplest case of a valuation-theoretic criterion going back
to Dumas \cite{Dumas}. We work throughout with the following family.

\begin{definition}\label{def:ED}
Let $p$ be a prime and let $m,n\ge1$ be coprime integers. A monic polynomial
$f=x^n+a_{n-1}x^{n-1}+\dots+a_0\in\Z[x]$ satisfies the \emph{Eisenstein--Dumas condition at $p$ of
slope $m/n$} if
\[
  \ord_p(a_0)=m
  \qquad\text{and}\qquad
  \ord_p(a_i)\ \ge\ \Big\lceil \frac{m(n-i)}{n}\Big\rceil \quad (1\le i\le n-1),
\]
where $\ord_p$ is the $p$-adic valuation. Equivalently, and without ceilings,
$m(n-i)\le n\,\ord_p(a_i)$ for all $0\le i\le n-1$, with equality $\ord_p(a_0)=m$; equivalently again,
the Newton polygon of $f$ with respect to $p$ is the single segment joining $(0,m)$ to $(n,0)$.
We write $\mathcal F_{p,m,n}$ for the set of monic $f \in \Z[x]$ of degree $n$ satisfying it.
\end{definition}

Taking $m=1$ gives back the $p$-Eisenstein polynomials, so $\mathcal F_{p,1,n}$ is the family
studied in \cite{HMT}. For $m>1$ the family is genuinely different: for instance
$x^3+25x^2+50x+50$ lies in $\mathcal F_{5,2,3}$ and is not Eisenstein at any prime.

Our first result is that the whole of the local theory survives the passage from $m=1$ to general
$m$; it is proved in Section~\ref{sec:dumas}.

\begin{theorem}\label{thm:main-ram}
Let $f \in \mathcal F_{p,m,n}$ and let $\theta$ be a root of $f$ in an algebraic closure of $\Q$,
$K=\Q(\theta)$. Let $\p$ be any prime of $\OK$ above $p$. Then
\[
  n\,\vp(\theta) \;=\; m\,\vp(p).
\]
Consequently $n\mid \vp(p)$, the degree $[K:\Q]$ equals $n$ (so that $f$ is irreducible), $\p$ is
the unique prime of $\OK$ above $p$, and
\[
  p\,\OK=\p^{\,n},\qquad \vp(\theta)=m .
\]
\end{theorem}

Thus one computation yields both the Eisenstein--Dumas irreducibility criterion and the total
ramification needed to run Heilbronn's criterion.

The second ingredient is the density of the family itself. Write
\[
  S(m,n)\;=\;\sum_{j=1}^{n}\Big\lceil \frac{mj}{n}\Big\rceil .
\]

\begin{theorem}\label{thm:main-density-family}
Let $p$ be a prime and let $m,n\ge1$ be coprime. Among the monic integer polynomials of degree $n$
of height at most $X$, the proportion lying in $\mathcal F_{p,m,n}$ tends, as $X\to\infty$, to
\[
  E_{p,m}(n)\;=\;\Big(1-\frac1p\Big)\,p^{-S(m,n)},
  \qquad\text{where}\qquad
  S(m,n)=\frac{(m+1)(n+1)}{2}-1 .
\]
For $m=1$ this is $p^{-n}-p^{-n-1}$, the Eisenstein density of \cite{Dubickas}.
\end{theorem}

The closed form for $S(m,n)$ is a pleasant consequence of the coprimality of $m$ and $n$; note
that $(m+1)(n+1)$ is always even in that case, since $m$ and $n$ are not both even.

Our main theorem is a single statement from which all the density bounds follow. For a prime $q$
let
\[
  C_q(n)\;=\;\frac{\#\{\text{monic } g\in\F_q[x] \text{ of degree } n \text{ with no root in }\F_q\}}{q^{\,n}} .
\]

\begin{theorem}\label{thm:master}
Let $n\ge2$, let $m\ge1$ be coprime to $n$, and let $p$ be a prime with $\gcd(p-1,n)=1$. Let $Q$
be a finite set of primes with $p\notin Q$, $\#Q\ge2$, such that
\begin{equation}\label{eq:pair-condition}
  \parbox{0.82\textwidth}{for all $q_1<q_2$ in $Q$ there are integers $u,v>0$ with
  $p=uq_1+vq_2$, $q_1\nmid u$ and $q_2\nmid v$.}
\end{equation}
Then, among the polynomials $f\in\mathcal F_{p,m,n}$ with $\Ht(f)\le X$, the proportion of those
for which $K_f$ is not norm-Euclidean satisfies
\[
  \liminf_{X\to\infty}\ \delta_{p,m,n}(X)
  \ \ge\
  \sum_{\substack{T\subseteq Q\\ \#T\ge2}}\ \prod_{q\in T}C_q(n)\prod_{q\in Q\setminus T}\bigl(1-C_q(n)\bigr).
\]
\end{theorem}

By Lemma~\ref{lem:frobenius} below, condition \eqref{eq:pair-condition} holds as soon as
$q_1^2q_2^2\le p$ for every pair $q_1<q_2$ in $Q$, and in particular whenever every element of $Q$
is at most $p^{1/4}$. The right-hand side of the displayed inequality is the exact density of the
subfamily of those $f$ having no root modulo $q$ for at least two primes $q\in Q$. Specialising $Q$ recovers, and interpolates between, the two
bounds of \cite{HMT}.

\begin{corollary}\label{cor:227}
Let $n\ge3$ and let $p\ge5$ be a prime with $\gcd(p-1,n)=1$, and let $m\ge1$ be coprime to $n$.
Then $\liminf_{X\to\infty}\delta_{p,m,n}(X)\ \ge\ 2/27$.
\end{corollary}

\begin{corollary}\label{cor:235}
Let $n \ge 3$, let $m \ge 1$ be coprime to $n$, and let $p\ge225$ be a prime with
$\gcd(p-1,n)=1$. Then
\[
  \liminf_{X\to\infty}\delta_{p,m,n}(X)\ \ge\ \frac{136}{675}\;=\;0.20148\ldots
\]
\end{corollary}

\begin{corollary}\label{cor:eps}
Let $n\ge3$, let $m\ge1$ be coprime to $n$, and let $p$ be a prime with $\gcd(p-1,n)=1$. Put
$t=\pi(\lfloor p^{1/4}\rfloor)$. Then
\[
  \liminf_{X\to\infty}\delta_{p,m,n}(X)\ \ge\ 1-\varepsilon(p),
  \qquad
  \varepsilon(p)=(1+t)\Big(\frac34\Big)^{t} .
\]
In particular $\varepsilon(p)\to0$ as $p\to\infty$, uniformly in $n$ and $m$.
\end{corollary}

Corollary~\ref{cor:235} illustrates the gain: for $p$ in the range where $t$ is still small,
$\varepsilon(p)>1$ and the bound of Corollary~\ref{cor:eps} is vacuous, while $2/27=0.074\ldots$
is far from optimal; a three-element $Q$ already gives $0.2014\ldots$.

Two remarks on the hypotheses. The condition $\gcd(p-1,n)=1$ forces $n$ to be odd, since $p-1$ is
even for $p$ odd; we therefore do not need to assume the parity of $n$, and, more importantly, we
do not need the step of \cite{HMT} that converts ``$-b$ is not a norm'' into ``$b$ is not a norm''
using $n$ odd. The reason is Lemma~\ref{lem:degree-one}, which is insensitive to signs.

Section~\ref{sec:dumas} proves Theorem~\ref{thm:main-ram}. Section~\ref{sec:heilbronn} contains
the non-norm lemma and an embedding-free proof of Heilbronn's criterion.
Section~\ref{sec:counting} gives the counting lemma, with one modulus per coefficient, which is
what the Eisenstein--Dumas condition requires. Section~\ref{sec:local} computes the local
densities, proving Theorem~\ref{thm:main-density-family} and the estimates for $C_q(n)$.
Section~\ref{sec:proof} proves Theorem~\ref{thm:master} and its corollaries. Section~\ref{sec:gcd}
weakens the condition $\gcd(p-1,n)=1$: Theorem~\ref{thm:gcd} there replaces it by the explicit
inequality
\[
  \bigl(\gcd(p-1,n)-1\bigr)\,p^{1/2}\bigl(1+\log p\bigr)\ <\ \frac{p}{q_1^{2}q_2^{2}}-2 ,
\]
in which $q_1,q_2$ are the two auxiliary primes. Since the bound is independent of $q_1$ and
$q_2$, the auxiliary primes may be taken as large as a fixed power of $p$, and
Corollary~\ref{cor:gcd-density} gives
$\liminf_{X\to\infty}\delta_{p,m,n}(X)\ge1-(1+\pi(Y))(3/4)^{\pi(Y)}$ for every $Y$ with
$8nY^{4}\le\lfloor p^{1/4}\rfloor$, for instance $Y=\lfloor p^{1/17}\rfloor$ once
$p\ge(8n)^{68}$.

\subsection*{Notation}
$K$ always denotes a number field, $\OK$ its ring of integers, $\N=\N_{K/\Q}$ the norm, and
$\vp$ the additive valuation attached to a prime $\p$ of $\OK$, normalised so that $\vp$ takes the
value $1$ at a uniformiser. For a prime $q$ of $\Z$ we also write $\ord_q$ for the $q$-adic valuation
on $\Z$. Letters $p,q$ always denote rational primes, and $\pi(Y)$ is the number of primes at
most $Y$.

\section{Eisenstein--Dumas polynomials and total ramification}\label{sec:dumas}

\begin{proof}[Proof of Theorem~\ref{thm:main-ram}]
Write $f=x^n+\sum_{i<n}a_ix^i$ and let $\theta$ be a root, $K=\Q(\theta)$, $d=[K:\Q]\le n$. As $f$
is monic, $\theta\in\OK$. Fix a prime $\p$ of $\OK$ above $p$ and set
\[
  w=\vp(\theta)\ \ge\ 0, \qquad e=\vp(p) .
\]
Since the ramification indices and residue degrees of the primes above $p$ satisfy
$\sum e_i f_i = d$, we have $e\le d\le n$. Note also $a_0\ne0$, because $\ord_p(a_0)=m$ is finite;
hence $\theta\ne0$, because $f(\theta)=0$ would otherwise give $a_0=0$.

Because $\vp$ restricted to $\Q$ is $e\,\ord_p$, we have $\vp(a_0)=em$ exactly, and, for every
$i<n$ with $a_i\ne0$,
\begin{equation}\label{eq:term-bound}
  n\,\vp\bigl(a_i\theta^i\bigr)\;=\;n\,e\,\ord_p(a_i)+n\,i\,w\;\ge\;e\,m(n-i)+n\,i\,w ,
\end{equation}
using $n\,\ord_p(a_i)\ge m(n-i)$. Terms with $a_i=0$ may be ignored throughout. We claim
\begin{equation}\label{eq:key}
  n\,w\;=\;e\,m .
\end{equation}

Suppose first that $nw>em$. Then for $1\le i\le n-1$, using \eqref{eq:term-bound},
\[
  n\,\vp\bigl(a_i\theta^i\bigr)\ \ge\ em(n-i)+niw\ >\ em(n-i)+i\,em\ =\ n\,em\ =\ n\,\vp(a_0),
\]
so $\vp(a_i\theta^i)>\vp(a_0)$ for every $i$ with $1\le i\le n-1$. By the ultrametric equality for
a sum with a unique term of least valuation,
$\vp\bigl(\sum_{i<n}a_i\theta^i\bigr)=\vp(a_0)=em$. But
$\theta^n=-\sum_{i<n}a_i\theta^i$, so $nw=em$, contradicting $nw>em$.

Suppose next that $nw<em$. Then $\vp(a_0)=em>nw$, and for $1\le i\le n-1$, since $n-i\ge1$,
\[
  n\,\vp\bigl(a_i\theta^i\bigr)\ \ge\ em(n-i)+niw\ >\ nw(n-i)+niw\ =\ n^2w ,
\]
so $\vp(a_i\theta^i)>nw$ for every $i<n$. Hence
$\vp\bigl(\sum_{i<n}a_i\theta^i\bigr)>nw$, while
$\vp(\theta^n)=nw$; this contradicts $\theta^n=-\sum_{i<n}a_i\theta^i$. This proves \eqref{eq:key}.

From \eqref{eq:key} we get $n\mid em$, and since $\gcd(m,n)=1$ this gives $n\mid e$. As
$e\le d\le n$ and $e\ge1$, we conclude $e=d=n$. In particular $[K:\Q]=n=\deg f$, so $f$ is the
minimal polynomial of $\theta$ and is irreducible. Moreover $\sum_i e_if_i=n$ with one summand
already equal to $e\cdot f(\p)\ge n$ forces $\p$ to be the only prime above $p$ and $f(\p)=1$;
therefore $p\OK=\p^{\,n}$. Finally $nw=em=nm$ gives $w=m$.
\end{proof}

\begin{corollary}[Eisenstein--Dumas criterion]\label{cor:dumas-irr}
A monic polynomial satisfying the Eisenstein--Dumas condition at some prime $p$ of slope $m/n$
with $\gcd(m,n)=1$ is irreducible over $\Q$.
\end{corollary}

\begin{remark}
For $m=1$ the theorem is the familiar statement that a $p$-Eisenstein polynomial is irreducible
and that $p$ is totally ramified in the field it generates; the proof above specialises to the
usual one. The point of the general case is that the extra room in the inequalities
$\ord_p(a_i)\ge \lceil m(n-i)/n \rceil$ costs nothing: only the two extreme cases $nw>em$ and
$nw<em$ have to be excluded, and both are excluded by comparing with the single term $a_0$.
\end{remark}

\section{Norms, degree-one primes, and Heilbronn's criterion}\label{sec:heilbronn}

Throughout this section $K$ is a number field of degree $n$ and $\theta\in\OK$ generates $K$.

\subsection{Which integers can be norms}

\begin{lemma}\label{lem:degree-one}
Let $\theta\in\OK$ be a root of a monic $f\in\Z[x]$ and let $q$ be a prime such that $f$ has no
root in $\F_q$. Then:
\begin{enumerate}
\item no prime $\q$ of $\OK$ above $q$ has residue degree $1$;
\item $\ord_q\bigl(\N(\alpha)\bigr)\ne1$ for every $\alpha\in\OK$;
\item if $u\in\Z$ and $q\nmid u$, then $uq$ is not the norm of an element of $\OK$.
\end{enumerate}
\end{lemma}

\begin{proof}
(1) If $\q$ lies above $q$ and has residue degree $1$, then $\OK/\q\cong\F_q$; reducing the
relation $f(\theta)=0$ modulo $\q$ exhibits the image of $\theta$ as a root of $f$ in $\F_q$, a
contradiction.

(2) Let $\alpha\in\OK$ be nonzero and factor $(\alpha)=\prod_i\q_i^{e_i}$. Taking absolute norms,
$|\N(\alpha)|=\prod_i \N(\q_i)^{e_i}$, and $\N(\q_i)$ is a power $q_i^{f_i}$ of the rational prime
$q_i$ below $\q_i$. Hence
\[
  \ord_q\bigl(\N(\alpha)\bigr)=\sum_{\q_i\mid q} f_i e_i .
\]
By (1) every $f_i$ occurring here is at least $2$, so the sum is either $0$ or at least $2$. For
$\alpha=0$ the norm is $0$ and the claim is vacuous.

(3) If $\N(\alpha)=uq$ with $q\nmid u$ then $\ord_q(\N(\alpha))=1$, contradicting (2).
\end{proof}

\begin{remark}\label{rem:lemma23}
Part (3) is the statement that the argument of \cite{HMT} requires. The lemma recorded there is
part (1) in the contrapositive form ``if $q$ is a norm from $\OK$ then $f$ has a root in $\F_q$'',
and in the application one needs to know not that $q_1$ is not a norm, but that $u q_1$ is not a
norm, where $p=uq_1+vq_2$; since the norms from $\OK$ are not stable under division this does not
formally follow. The remedy is exactly part (2): the $q$-adic valuation of a norm is never $1$.
This also explains the role of the conditions $q_1\nmid u$ and $q_2\nmid v$ in
Lemma~\ref{lem:frobenius} below. Finally, since (2) makes no reference to signs, $-vq_2$ is
excluded on the same footing as $uq_1$, and there is no need to assume that $n$ is odd in order to
pass between $b$ and $-b$.
\end{remark}

\subsection{The norm modulo a totally ramified prime}

The following lemma is the heart of the matter. It replaces the classical argument through the
conjugates $\rho^\sigma$ of an element, which requires a Galois closure to be meaningful, by a
statement about a determinant.

\begin{lemma}\label{lem:norm-power}
Let $A$ be a commutative ring, let $B$ be a free $A$-algebra of finite rank $n$, and let $p\in A$
be such that $A/(p)$ is reduced. Let $I\subseteq B$ be an ideal with $I^{\,n}\subseteq pB$. If
$\rho\in B$ and $x\in A$ satisfy $\rho-x\in I$, then
\[
  \N_{B/A}(\rho)\ \equiv\ x^{\,n} \pmod{p}.
\]
\end{lemma}

\begin{proof}
Fix an $A$-basis of $B$ and let $M_y$ denote the matrix of multiplication by $y\in B$ in that
basis, so that $\N_{B/A}(y)=\det M_y$. Put $t=\rho-x\in I$, so $M_\rho=x\,\mathrm{Id}+M_t$.
Reduce all matrix entries modulo $p$ and write $\overline{\phantom{x}}$ for the images; since the
determinant is a polynomial in the entries, $\overline{\det M_\rho}=\det\overline{M_\rho}$.

Because $t\in I$ we have $t^{\,n}\in I^{\,n}\subseteq pB$, say $t^{\,n}=p\,y$; as
$y\mapsto M_y$ is a ring homomorphism which is $A$-linear, $M_t^{\,n}=M_{t^n}=p\,M_y$ has all its
entries in $(p)$, and therefore $\overline{M_t}^{\,n}=0$: the reduced matrix $\overline{M_t}$ is
nilpotent. Over any commutative ring the characteristic polynomial of a nilpotent matrix differs
from $X^n$ by a polynomial with nilpotent coefficients; since $A/(p)$ is reduced, the
characteristic polynomial of $-\overline{M_t}$ is exactly $X^n$. Evaluating it at $\bar x$ gives
\[
  \det\bigl(\bar x\,\mathrm{Id}+\overline{M_t}\bigr)
  =\det\bigl(\bar x\,\mathrm{Id}-(-\overline{M_t})\bigr)=\bar x^{\,n},
\]
which is the assertion.
\end{proof}

\begin{corollary}\label{cor:norm-power-K}
Let $K$ be a number field of degree $n$ and let $\p$ be a prime of $\OK$ with $\p^{\,n}=p\OK$ for
a rational prime $p$. Then for every $\rho\in\OK$ and every $x\in\Z$ with $\rho\equiv x \pmod \p$,
\[
  \N(\rho)\ \equiv\ x^{\,n}\pmod p .
\]
\end{corollary}

\begin{proof}
Apply Lemma~\ref{lem:norm-power} with $A=\Z$, $B=\OK$, which is free of rank $n$ over $\Z$, with
$I=\p$, and with $\Z/(p)=\F_p$, which is reduced. The hypothesis $I^n\subseteq pB$ is the
assumption $\p^n=p\OK$.
\end{proof}

Since $\p^{\,n}=p\OK$ forces the residue field $\OK/\p$ to be $\F_p$, every $\rho\in\OK$ is
congruent modulo $\p$ to some rational integer $x$, and Corollary~\ref{cor:norm-power-K} says
precisely that the norm map induces $y\mapsto y^{\,n}$ on residue fields.

\subsection{Proof of Heilbronn's criterion}

\begin{lemma}\label{lem:principal}
Suppose $K$ is norm-Euclidean. Then every ideal of $\OK$ is principal.
\end{lemma}

\begin{proof}
Let $I\ne 0$ be an ideal and choose $\pi\in I\setminus\{0\}$ with $|\N(\pi)|$ minimal among the
nonzero elements of $I$; this is possible since $|\N|$ takes values in the nonnegative integers.
For $\alpha\in I$ write $\alpha=\gamma\pi+\rho$ with $|\N(\rho)|<|\N(\pi)|$. Then
$\rho=\alpha-\gamma\pi\in I$, so $\rho=0$ by minimality, and $\alpha\in(\pi)$. Hence $I=(\pi)$.
\end{proof}

\begin{proof}[Proof of Theorem~\ref{thm:heilbronn-intro}]
Let $\p$ be the prime above $p$, so that $p\OK=\p^{\,n}$, and suppose for contradiction that $K$
is norm-Euclidean. By Lemma~\ref{lem:principal}, $\p=(\pi)$ for some $\pi\in\OK$. Taking absolute
norms in $\p^{\,n}=p\OK$ gives $\N(\p)^n=p^n$, so $\N(\p)=p$, that is $|\N(\pi)|=p$.

Since $a$ is an $n$-th power residue modulo $p$, choose $x\in\Z$ with $x^n\equiv a \pmod p$.
Dividing $x$ by $\pi$ produces $\gamma,\rho\in\OK$ with
\[
  x=\gamma\pi+\rho,\qquad |\N(\rho)|<|\N(\pi)|=p .
\]
Then $\rho\equiv x \pmod{\p}$, so Corollary~\ref{cor:norm-power-K} gives
$\N(\rho)\equiv x^n\equiv a \pmod p$. Now $0<a<p$ and $-p<\N(\rho)<p$, so
$\N(\rho)-a\in(-2p,p)$ is divisible by $p$ and therefore equals $0$ or $-p$; that is,
$\N(\rho)=a$ or $\N(\rho)=a-p=-b$. Either way one of $a$, $-b$ is a norm, contrary to hypothesis.
\end{proof}

\subsection{Representing $p$ by two primes}

\begin{lemma}\label{lem:frobenius}
Let $q_1<q_2$ be primes and let $p\ge q_1^2q_2^2$. Then there are positive integers $u,v$ with
\[
  p=uq_1+vq_2,\qquad q_1\nmid u,\qquad q_2\nmid v .
\]
\end{lemma}

\begin{proof}
Fix $r_1\in\{1,\dots,q_1-1\}$ and $r_2\in\{1,\dots,q_2-1\}$. From $p\ge q_1^2q_2^2$,
\[
  p-r_1q_1-r_2q_2\ \ge\ q_1^2q_2^2-(q_1-1)q_1-(q_2-1)q_2\ \ge\ (q_1^2-1)(q_2^2-1).
\]
Since $q_1^2$ and $q_2^2$ are coprime, every integer at least $(q_1^2-1)(q_2^2-1)$ is a
nonnegative integral combination of $q_1^2$ and $q_2^2$ (the Frobenius coin problem), so there are
$t_1,t_2\ge0$ with $p-r_1q_1-r_2q_2=q_1^2t_1+q_2^2t_2$. Then
\[
  p=q_1(q_1t_1+r_1)+q_2(q_2t_2+r_2)=uq_1+vq_2
\]
with $u=q_1t_1+r_1$ and $v=q_2t_2+r_2$ positive and satisfying $q_1\nmid u$, $q_2\nmid v$.
\end{proof}

\begin{lemma}\label{lem:heilbronn-applies}
Let $n\ge1$ and let $p$ be a prime with $\gcd(p-1,n)=1$. Let $K=\Q(\theta)$ with $\theta\in\OK$ a
root of a monic $f\in\Z[x]$ of degree $n$, and suppose $p$ is totally ramified in $K$. Let
$q_1<q_2$ be primes different from $p$ such that $f$ has no root modulo $q_1$ and none modulo
$q_2$, and suppose $p=uq_1+vq_2$ with $u,v>0$, $q_1\nmid u$ and $q_2\nmid v$. Then $K$ is not
norm-Euclidean.
\end{lemma}

\begin{proof}
Put $a=uq_1$ and $b=vq_2$, so $p=a+b$ with $a,b>0$. Since $\gcd(n,p-1)=1$, the map $y\mapsto y^n$
is a bijection of $\F_p$: it is injective on the cyclic group $\F_p^\times$ of order $p-1$ because
$\gcd(n,p-1)=1$, and fixes $0$. Hence every integer, in particular $a$, is an $n$-th power residue
modulo $p$. By Lemma~\ref{lem:degree-one}(3) applied to $q_1$, the integer $a=uq_1$ is not a norm;
applied to $q_2$, and using that the statement of that lemma is about the valuation of the norm
and hence indifferent to sign, the integer $-b=-vq_2$ is not a norm either. Theorem
\ref{thm:heilbronn-intro} now applies.
\end{proof}

\section{Counting polynomials with local conditions}\label{sec:counting}

The Eisenstein--Dumas condition prescribes the coefficient $a_i$ modulo a power of $p$ that
\emph{depends on $i$}. The following counting lemma is therefore stated with one modulus per
coordinate; this costs nothing and makes the application immediate.

\begin{lemma}\label{lem:box}
Let $\iota$ be a finite set of size $N$, let $M:\iota\to\Z_{\ge1}$, and let
$S\subseteq\prod_{i\in\iota}\Z/M_i\Z$. For $X\in\Z_{\ge 0}$ put
\[
  \mathcal N(X)=\#\bigl\{a\in\Z^{\iota}\ :\ |a_i|\le X \text{ for all } i,\ (a_i \bmod M_i)_{i}\in S\bigr\}.
\]
Then
\[
  \mathcal N(X)\ \le\ \#S\cdot\prod_{i\in\iota}\Bigl(\frac{2X+1}{M_i}+1\Bigr),
\]
and, provided $M_i\le 2X+1$ for every $i$,
\[
  \mathcal N(X)\ \ge\ \#S\cdot\prod_{i\in\iota}\Bigl(\frac{2X+1}{M_i}-1\Bigr).
\]
Consequently
\[
  \lim_{X\to\infty}\frac{\mathcal N(X)}{(2X+1)^N}\;=\;\frac{\#S}{\prod_{i\in\iota}M_i}.
\]
\end{lemma}

\begin{proof}
Sorting the tuples $a$ by their reduction gives
\[
  \mathcal N(X)=\sum_{s\in S}\ \prod_{i\in\iota}\#\{x\in\Z: |x|\le X,\ x\equiv s_i \bmod M_i\},
\]
because the conditions on the separate coordinates are independent. The interval $[-X,X]$ contains
$2X+1$ integers, and the integers in it congruent to a fixed residue modulo $M$ are equally spaced
with gap $M$, so their number lies between $(2X+1)/M-1$ and $(2X+1)/M+1$. The upper bound follows
at once. For the lower bound one needs the factors $(2X+1)/M_i-1$ to be nonnegative before
multiplying the inequalities, which is exactly the hypothesis $M_i\le 2X+1$. Finally, dividing by
$(2X+1)^N$ turns the two bounds into
$\#S\prod_i\bigl(M_i^{-1}\mp(2X+1)^{-1}\bigr)$, and both tend to $\#S/\prod_iM_i$.
\end{proof}

\begin{remark}
The hypothesis $M_i\le 2X+1$ in the lower bound cannot be dropped: without it an even number of
factors $(2X+1)/M_i-1$ may be negative and their product positive, and the resulting inequality is
false. Since we let $X\to\infty$ with the moduli fixed, the hypothesis is harmless.
\end{remark}

By the Chinese Remainder Theorem, conditions imposed modulo powers of distinct primes are
independent, and a set $S$ of product form $\prod_i T_i$ has $\#S/\prod_i M_i=\prod_i(\#T_i/M_i)$.
We record the form in which the lemma will be used.

\begin{proposition}\label{prop:localdensities}
Let $p$ be a prime, let $m,n\ge1$ be coprime, and let $Q$ be a finite set of primes not containing
$p$. For each $q\in Q$ let $R_q\subseteq(\Z/q\Z)^n$ be a set of coefficient vectors. Then the
number $N^{*}(X)$ of monic $f\in\Z[x]$ of degree $n$ with $\Ht(f)\le X$ which lie in
$\mathcal F_{p,m,n}$ and whose coefficient vector lies in $R_q$ modulo $q$ for every $q\in Q$
satisfies
\[
  \lim_{X\to\infty}\frac{N^{*}(X)}{(2X+1)^n}\;=\;E_{p,m}(n)\cdot\prod_{q\in Q}\frac{\#R_q}{q^{\,n}},
\]
where $E_{p,m}(n)$ is the density of Theorem~\ref{thm:main-density-family}.
\end{proposition}

\begin{proof}
The conditions defining $\mathcal F_{p,m,n}$ are conditions on $a_i$ modulo $p^{m+1}$ for $i=0$
and modulo $p^{\lceil m(n-i)/n\rceil}$ for $1\le i\le n-1$; the conditions at $q\in Q$ are
conditions modulo $q$. Apply Lemma~\ref{lem:box} with $M_i$ the product of the relevant prime
powers at each coordinate and with $S$ the set cut out by all the conditions, which by the Chinese
Remainder Theorem has $\#S/\prod_iM_i$ equal to the product of the corresponding local densities.
The local density at $p$ is computed in Proposition~\ref{prop:E} below.
\end{proof}

\section{Local densities}\label{sec:local}

\subsection{The Eisenstein--Dumas density}

\begin{proposition}\label{prop:E}
Let $p$ be a prime and let $m,n\ge1$ be coprime. The density of $\mathcal F_{p,m,n}$ among monic
integer polynomials of degree $n$ is
\[
  E_{p,m}(n)=\Bigl(1-\frac1p\Bigr)p^{-S(m,n)},\qquad S(m,n)=\sum_{j=1}^{n}\Bigl\lceil\frac{mj}{n}\Bigr\rceil
  =\frac{(m+1)(n+1)}{2}-1 .
\]
\end{proposition}

\begin{proof}
Set $c_i=\lceil m(n-i)/n\rceil$ for $0\le i\le n-1$, so that $c_0=m$. The conditions are
$\ord_p(a_0)=m$, which is a condition modulo $p^{m+1}$ satisfied by exactly
$p^{m+1}(p^{-m}-p^{-m-1})$ residues, and $\ord_p(a_i)\ge c_i$ for $1\le i\le n-1$, which is satisfied
by a proportion $p^{-c_i}$ of residues. The conditions on distinct coefficients are independent,
so by Lemma~\ref{lem:box} the density exists and equals
\[
  \Bigl(\frac1{p^{m}}-\frac1{p^{m+1}}\Bigr)\prod_{i=1}^{n-1}p^{-c_i}
  =\Bigl(1-\frac1p\Bigr)p^{-\left(m+\sum_{i=1}^{n-1}c_i\right)} .
\]
Substituting $j=n-i$ turns $m+\sum_{i=1}^{n-1}c_i$ into $\sum_{j=1}^{n}\lceil mj/n\rceil=S(m,n)$.

For the closed form, note that $\gcd(m,n)=1$ implies $n\nmid mj$ for $1\le j\le n-1$, so
$\lceil mj/n\rceil=\lfloor mj/n\rfloor+1$ for those $j$, while the term $j=n$ contributes $m$.
Pairing $j$ with $n-j$ and using $\lfloor y\rfloor+\lfloor -y\rfloor=-1$ for $y\notin\Z$,
\[
  \Bigl\lfloor\frac{mj}{n}\Bigr\rfloor+\Bigl\lfloor\frac{m(n-j)}{n}\Bigr\rfloor=m-1
  \qquad(1\le j\le n-1),
\]
so that $2\sum_{j=1}^{n-1}\lfloor mj/n\rfloor=(m-1)(n-1)$. Hence
\[
  S(m,n)=m+(n-1)+\frac{(m-1)(n-1)}{2}=\frac{(m+1)(n+1)}{2}-1 . \qedhere
\]
\end{proof}

Note that $m$ and $n$ coprime are not both even, so $(m-1)(n-1)$ is even and $S(m,n)$ is an
integer, as it must be. For $m=1$ we get $S(1,n)=n$ and
$E_{p,1}(n)=p^{-n}-p^{-n-1}$, the density of $p$-Eisenstein polynomials computed by Dubickas
\cite{Dubickas}.

\subsection{Monic polynomials over a finite field without roots}

\begin{proposition}\label{prop:rootless}
Let $\F$ be a finite field with $q$ elements, let $n\ge 1$, and let $T\subseteq\F$. The number of
monic polynomials of degree $n$ over $\F$ with no root in $T$ is
\[
  A_q(n,T)=\sum_{k=0}^{n}(-1)^k\binom{\#T}{k}q^{\,n-k} .
\]
In particular, writing $A_q(n)=A_q(n,\F)$ and $C_q(n)=A_q(n)/q^{\,n}$, we have
\[
  A_q(n)=\sum_{k=0}^n(-1)^k\binom{q}{k}q^{\,n-k},
  \qquad\text{and}\qquad
  C_q(n)=\Bigl(1-\frac1q\Bigr)^{q}\ \text{ if } n\ge q .
\]
\end{proposition}

\begin{proof}
For $r\in T$ let $S_r$ be the set of monic polynomials of degree $n$ vanishing at $r$. By
inclusion--exclusion, the number of monic polynomials of degree $n$ avoiding every root in $T$ is
$\sum_{U\subseteq T}(-1)^{\#U}\#\bigl(\bigcap_{r\in U}S_r\bigr)$. Fix $U$ with $\#U=k$ and put
$g_U=\prod_{r\in U}(x-r)$, a monic polynomial of degree $k$. Since the elements of $U$ are
distinct, a polynomial vanishes on $U$ if and only if it is divisible by $g_U$; and
$h\mapsto g_Uh$ is a bijection from the monic polynomials of degree $n-k$ onto the monic
polynomials of degree $n$ divisible by $g_U$, with inverse $f\mapsto f/g_U$. Hence
$\#\bigcap_{r\in U}S_r=q^{\,n-k}$ when $k\le n$. When $k>n$ the intersection is empty, because a
nonzero polynomial of degree $n$ has at most $n$ roots. Grouping the subsets $U$ by their size
gives the stated formula, the terms with $k>n$ being absent and the terms with $k>\#T$ vanishing
because $\binom{\#T}{k}=0$.

For $T=\F$ and $n\ge q$ the binomial coefficients $\binom qk$ vanish for $k>q$, so
\[
  A_q(n)=\sum_{k=0}^{q}(-1)^k\binom qk q^{\,n-k}
  =q^{\,n-q}\sum_{k=0}^{q}\binom qk(-1)^kq^{\,q-k}
  =q^{\,n-q}(q-1)^q,
\]
by the binomial theorem, whence $C_q(n)=(1-1/q)^q$.
\end{proof}

\begin{proposition}\label{prop:bounds}
For $n\ge2$ and any prime power $q\ge2$,
\[
  \frac14\ \le\ \frac{q^2-1}{3q^2}\ \le\ C_q(n)\ \le\ \frac{q-1}{2q}\ <\ \frac12 .
\]
\end{proposition}

\begin{proof}
Dividing the formula of Proposition~\ref{prop:rootless} by $q^n$ gives
$C_q(n)=\sum_{k=0}^{n}(-1)^k\varphi(k)$ with $\varphi(k)=\binom qk q^{-k}$. We have
$\varphi(0)=\varphi(1)=1$, so the first two terms cancel and
\[
  C_q(n)=\sum_{k=2}^{n}(-1)^k\varphi(k).
\]
The sequence $\varphi$ is nonnegative and nonincreasing. Indeed
$(k+1)\binom{q}{k+1}=(q-k)\binom qk\le q\binom qk$, and $k+1\ge1$, so
$\binom{q}{k+1}\le q\binom qk$; dividing by $q^{k+1}$ gives
$\varphi(k+1)\le\varphi(k)$.

We use the following elementary bound: if $\psi$ is nonincreasing and nonnegative then, for
every $d\ge0$,
\[
  0\ \le\ \sum_{k=0}^{d-1}(-1)^k\psi(k)\ \le\ \psi(0),
\]
the empty sum being $0$. This is immediate by induction on $d$: for $d\ge1$,
\[
  \sum_{k=0}^{d-1}(-1)^k\psi(k)=\psi(0)-\sum_{k=0}^{d-2}(-1)^k\psi(k+1),
\]
and by the inductive hypothesis applied to the shifted sequence $k\mapsto\psi(k+1)$ the
subtracted sum lies in $[0,\psi(1)]\subseteq[0,\psi(0)]$. Applying this with $\psi(k)=\varphi(k+2)$ gives $C_q(n)\le\varphi(2)$, and applying
it once more after peeling off the first term,
$C_q(n)=\varphi(2)-\sum_{k=0}^{n-3}(-1)^k\varphi(k+3)\ge\varphi(2)-\varphi(3)$.

Finally $\varphi(2)=\binom q2/q^2=(q-1)/(2q)$ and $\varphi(3)=\binom q3/q^3=(q-1)(q-2)/(6q^2)$, so
\[
  \varphi(2)-\varphi(3)=\frac{(q-1)\bigl(3q-(q-2)\bigr)}{6q^2}=\frac{(q-1)(q+1)}{3q^2}=\frac{q^2-1}{3q^2},
\]
and $q\ge 2$ gives $(q^2-1)/(3q^2)\ge 1/4$ and $(q-1)/(2q)<1/2$.
\end{proof}

\begin{remark}
The bounds of Proposition~\ref{prop:bounds} appear in \cite{HMT}. The argument there deduces that
$\varphi$ is decreasing from the computation $\varphi(k)/\varphi(k+1)=(k+1)q/(q-k)>0$; positivity
of the ratio does not imply monotonicity, and in fact $\varphi$ is not strictly decreasing, since
$\varphi(0)=\varphi(1)=1$. The correct statement is that the ratio is at least $1$, with strict
inequality for $k\ge1$, which is what the proof above uses; the conclusion is unaffected.
\end{remark}

For the two smallest primes, Proposition~\ref{prop:rootless} gives the exact values
\[
  C_2(n)=\Bigl(1-\tfrac12\Bigr)^2=\tfrac14\ (n\ge2),\qquad
  C_3(n)=\Bigl(1-\tfrac13\Bigr)^3=\tfrac{8}{27}\ (n\ge3),
\]
while Proposition~\ref{prop:bounds} gives $C_5(n)\ge 8/25$ for all $n\ge2$.

\section{Proof of the main theorem}\label{sec:proof}

\begin{proof}[Proof of Theorem~\ref{thm:master}]
Fix $n$, $m$, $p$ and $Q$ as in the statement. For a monic $f$ of degree $n$ let
\[
  R(f)=\{q\in Q:\ f \text{ has no root modulo } q\}.
\]
We claim that if $\#R(f)\ge2$ and $f\in\mathcal F_{p,m,n}$ then $K_f$ is not norm-Euclidean.
Indeed, choose $q_1<q_2$ in $R(f)$. By Theorem~\ref{thm:main-ram}, $f$ is irreducible, so
$K_f=\Q(\theta)$ with $\theta$ a root of $f$, and $p$ is totally ramified in $K_f$; note
$\theta\in\OK$ since $f$ is monic. By \eqref{eq:pair-condition} we may write
$p=uq_1+vq_2$ with $u,v>0$, $q_1\nmid u$, $q_2\nmid v$. Now Lemma~\ref{lem:heilbronn-applies}
applies and $K_f$ is not norm-Euclidean.

It remains to compute the density of $\{f\in\mathcal F_{p,m,n}: \#R(f)\ge2\}$ inside
$\mathcal F_{p,m,n}$. Whether $q\in R(f)$ depends only on the reduction of the coefficient vector
of $f$ modulo $q$, and by Proposition~\ref{prop:rootless} the proportion of coefficient vectors in
$(\Z/q\Z)^n$ for which the associated monic polynomial has no root modulo $q$ is exactly $C_q(n)$.
For a subset $T\subseteq Q$ let $R_q=\{$vectors with no root mod $q\}$ for $q\in T$ and its
complement for $q\in Q\setminus T$; then Proposition~\ref{prop:localdensities} shows that the
number of $f\in\mathcal F_{p,m,n}$ with $\Ht(f)\le X$ and $R(f)=T$, divided by $(2X+1)^n$, tends to
\[
  E_{p,m}(n)\prod_{q\in T}C_q(n)\prod_{q\in Q\setminus T}\bigl(1-C_q(n)\bigr).
\]
Taking $Q=\emptyset$ in Proposition~\ref{prop:localdensities} shows that the denominator
$N_{p,m,n}(X)=\#\{f\in\mathcal F_{p,m,n}:\Ht(f)\le X\}$ satisfies
$N_{p,m,n}(X)/(2X+1)^n\to E_{p,m}(n)>0$. Summing the finitely many disjoint classes with $\#T\ge2$
and dividing, the proportion of $f\in\mathcal F_{p,m,n}$ with $\Ht(f)\le X$ and $\#R(f)\ge2$ tends
to $\sum_{\#T\ge2}\prod_{q\in T}C_q(n)\prod_{q\notin T}(1-C_q(n))$. Since every such $f$ generates
a field that is not norm-Euclidean, the lower density of the latter family is at least this
number.
\end{proof}

It is convenient to record the complementary form of the bound. Writing $c_q=C_q(n)$,
\begin{equation}\label{eq:complement}
  \sum_{\substack{T\subseteq Q\\ \#T\ge2}}\prod_{q\in T}c_q\prod_{q\notin T}(1-c_q)
  \;=\;1-\prod_{q\in Q}(1-c_q)\Bigl(1+\sum_{q\in Q}\frac{c_q}{1-c_q}\Bigr),
\end{equation}
since the omitted classes are $T=\emptyset$ and the singletons.

The density computation in the proof above did not use the hypothesis
\eqref{eq:pair-condition}, which entered only when passing from ``two rootless primes'' to
``not norm-Euclidean''. We isolate it for use in Section~\ref{sec:gcd}.

\begin{proposition}\label{prop:two-rootless}
Let $n\ge2$, let $m\ge1$ be coprime to $n$, let $p$ be a prime and let $Q$ be a finite set of
primes with $p\notin Q$. Among the $f\in\mathcal F_{p,m,n}$ with $\Ht(f)\le X$, the proportion of
those having no root modulo $q$ for at least two $q\in Q$ tends, as $X\to\infty$, to
$\sum_{\#T\ge2}\prod_{q\in T}C_q(n)\prod_{q\notin T}(1-C_q(n))$, and this quantity is at least
$1-(1+\#Q)(3/4)^{\#Q}$.
\end{proposition}

\begin{proof}
The limit is the content of the second half of the proof of Theorem~\ref{thm:master}. For the
lower bound use \eqref{eq:complement} together with $C_q(n)\ge\frac14$ and $C_q(n)<\frac12$,
which hold by Proposition~\ref{prop:bounds} because $n\ge2$: the first gives
$1-C_q(n)\le\frac34$ and the second gives $C_q(n)/(1-C_q(n))\le1$.
\end{proof}

\begin{proof}[Proof of Corollary~\ref{cor:227}]
Suppose first that $p\notin\{7,11,19\}$ and take $Q=\{2,3\}$. If $p\ge36=2^23^2$ then
\eqref{eq:pair-condition} holds by Lemma~\ref{lem:frobenius}. If $p<36$ then
$p\in\{5,13,17,23,29,31\}$, and \eqref{eq:pair-condition} is verified directly by
\[
  5=2\cdot1+3\cdot1,\quad 13=2\cdot 5+3\cdot 1,\quad 17=2\cdot7+3\cdot1,
\]
\[
  23=2\cdot1+3\cdot7,\quad 29=2\cdot13+3\cdot1,\quad 31=2\cdot5+3\cdot7,
\]
in each of which the coefficient of $2$ is odd and the coefficient of $3$ is prime to $3$. With
$\#Q=2$ the sum in Theorem~\ref{thm:master} has the single term $C_2(n)C_3(n)$, and since $n\ge3$
Proposition~\ref{prop:rootless} gives
$C_2(n)C_3(n)=\frac14\cdot\frac8{27}=\frac{2}{27}$.

If $p\in\{7,11,19\}$ take instead $Q=\{2,5\}$ and use
\[
  7=2\cdot1+5\cdot1,\qquad 11=2\cdot3+5\cdot1,\qquad 19=2\cdot7+5\cdot1 ,
\]
which again satisfy \eqref{eq:pair-condition}. The bound is then
$C_2(n)C_5(n)\ge\frac14\cdot\frac8{25}=\frac2{25}>\frac2{27}$ by Propositions~\ref{prop:rootless}
and~\ref{prop:bounds}.
\end{proof}

\begin{proof}[Proof of Corollary~\ref{cor:235}]
Take $Q=\{2,3,5\}$; the largest pair product is $3^25^2=225\le p$, so
\eqref{eq:pair-condition} holds by Lemma~\ref{lem:frobenius}. The left-hand side of
\eqref{eq:complement} is the probability that at least two of three independent events of
probabilities $c_2,c_3,c_5$ occur; that event is increasing in each coordinate, so the quantity is
nondecreasing in each $c_q$. With $c_2=\frac14$, $c_3=\frac8{27}$ and $c_5\ge\frac8{25}$ it is
therefore at least its value at $(\frac14,\frac8{27},\frac8{25})$, namely
\[
  1-\frac{969+323+408+456}{2700}=1-\frac{2156}{2700}=\frac{544}{2700}=\frac{136}{675}. \qedhere
\]
\end{proof}

\begin{proof}[Proof of Corollary~\ref{cor:eps}]
Let $Y=\lfloor p^{1/4}\rfloor$ and let $Q$ be the set of primes $q\le Y$, so $\#Q=t$. If
$t\le1$ then $\varepsilon(p)\ge1$ and there is nothing to prove. If $t\ge2$, then for $q_1<q_2$
in $Q$ we have $q_1^2q_2^2\le Y^4\le p$, so \eqref{eq:pair-condition} holds by
Lemma~\ref{lem:frobenius}, and $p\notin Q$.
Theorem~\ref{thm:master} together with the lower bound of Proposition~\ref{prop:two-rootless}
gives $\liminf\delta_{p,m,n}(X)\ge1-(1+t)(3/4)^t$. Finally $Y\to\infty$ as $p\to\infty$, so
$t=\pi(Y)\to\infty$ and $(1+t)(3/4)^t\to0$; the bound involves neither $n$ nor $m$.
\end{proof}

Since $\mathcal F_{p,m,n}$ has positive density inside the monic polynomials of degree $n$
(Theorem~\ref{thm:main-density-family}), we obtain in particular:

\begin{corollary}\label{cor:positive}
Let $n\ge3$ be odd and let $m\ge1$ be coprime to $n$. A positive proportion of the monic integer
polynomials of degree $n$ satisfying the Eisenstein--Dumas condition of slope $m/n$ at $5$
generate fields that are not norm-Euclidean; the same holds inside the family of all monic
polynomials of degree $n$.
\end{corollary}

\begin{proof}
Apply Corollary~\ref{cor:227} with $p=5$, which satisfies $\gcd(4,n)=1$ for odd $n$, and combine
with Theorem~\ref{thm:main-density-family}.
\end{proof}

\section{Weakening the coprimality condition}\label{sec:gcd}

The hypothesis $\gcd(p-1,n)=1$ was used only to guarantee that $a=uq_1$ is an $n$-th power residue
modulo $p$. For a fixed $n$ there are infinitely many primes $p$ with $\gcd(p-1,n)=1$, but the
condition fails identically when $n$ is even. It can be relaxed at the cost of taking $p$ large,
by choosing $u$ inside a suitable arithmetic progression; the price is a character sum estimate,
for which we use the classical inequality of P\'olya and Vinogradov \cite{Polya,Vinogradov} in the
following explicit form, for a prime modulus. For a non-principal Dirichlet character $\chi$
modulo a prime $p$, any $c\in\Z/p\Z$ and any $N\ge0$,
\begin{equation}\label{eq:PV}
  \Bigl|\sum_{t<N}\chi(c+t)\Bigr|\ \le\ p^{1/2}\bigl(1+\log p\bigr).
\end{equation}
Indeed, a non-principal character modulo a prime is primitive, so its finite Fourier expansion has
coefficients $\overline\chi(-k)\tau(\chi)$ with $|\tau(\chi)|=p^{1/2}$; the resulting complete
exponential sums are geometric series bounded by $1/|\sin(\pi k/p)|$; and
$\sin(\pi a/p)\ge 2\min(a,p-a)/p$ by concavity, so that
\[
  \sum_{a=1}^{p-1}\frac1{\sin(\pi a/p)}
  \ \le\ \frac p2\sum_{a=1}^{p-1}\Bigl(\frac1a+\frac1{p-a}\Bigr)
  \ =\ p\,H_{p-1}\ \le\ p\,(1+\log p),
\]
which after division by $p$ and multiplication by $|\tau(\chi)|=p^{1/2}$ gives \eqref{eq:PV}.

The point of the argument below is that the three conditions to be imposed on $u$ --- two
congruence conditions and one multiplicative condition --- can be reduced to a
\emph{single} arithmetic progression together with the multiplicative condition. Consequently only
characters modulo $p$ occur, no reduction of imprimitive characters to primitive ones is needed,
and the resulting criterion is completely explicit.

\begin{theorem}\label{thm:gcd}
Let $n\ge2$, let $m\ge1$ be coprime to $n$, let $p$ be a prime and put $g=\gcd(p-1,n)$. Let
$q_1\ne q_2$ be primes different from $p$ and suppose that
\begin{equation}\label{eq:gcd-condition}
  (g-1)\,p^{1/2}\bigl(1+\log p\bigr)\ <\ \frac{p}{q_1^{2}q_2^{2}}-2 .
\end{equation}
If $f\in\mathcal F_{p,m,n}$ has no root modulo $q_1$ and no root modulo $q_2$, then $K_f$ is not
norm-Euclidean.
\end{theorem}

\begin{proof}
By Theorem~\ref{thm:main-ram}, $f$ is irreducible and $p$ is totally ramified in $K_f$, so by
Lemma~\ref{lem:heilbronn-applies} it suffices to produce positive integers $u,v$ with
\begin{equation}\label{eq:target}
  p=uq_1+vq_2,\qquad q_1\nmid u,\qquad q_2\nmid v,\qquad uq_1 \text{ an } n\text{-th power residue mod } p .
\end{equation}

Let $\psi$ be a Dirichlet character modulo $p$ with $\psi^{g}=1$ whose kernel is exactly the group
of $n$-th powers. Such a $\psi$ exists: the group $(\Z/p\Z)^{\times}$ is cyclic of order $N=p-1$,
so it is isomorphic to the group $\mu_N\subset\mathbb{C}^{\times}$ of $N$-th roots of unity; composing
such an isomorphism with the inclusion gives an injective character $\phi$, and
$\psi=\phi^{N/g}$ satisfies $\psi^{g}=\phi^{N}=1$ while $\psi(y)=1$ forces $y^{N/g}=1$, which for
a cyclic group of order $N$ says exactly that $y$ is an $n$-th power, the $n$-th powers being the
unique subgroup of index $g=\gcd(N,n)$. Thus for $p\nmid y$ we have $\psi(y)=1$ if and only if $y$
is an $n$-th power residue modulo $p$. Moreover $\psi$ has order exactly $g$: since $\phi$ is
injective, $\psi^{k}=\phi^{kN/g}=1$ holds if and only if $N\mid kN/g$, that is if and only if
$g\mid k$; in particular $\psi^{k}\ne1$ for $1\le k\le g-1$.

Instead of first choosing $u$ and then correcting for a possible divisibility $q_2\mid v$, we
impose the condition $q_2\nmid v$ from the start, as a congruence modulo $q_2^2$. Indeed, if
$p=uq_1+vq_2$ then
\[
  q_2\mid v \iff q_2^{2}\mid p-uq_1 \iff uq_1\equiv p \pmod{q_2^{2}} .
\]
Choose $c$ modulo $q_2^{2}$ with $cq_1\equiv p+q_1q_2 \pmod{q_2^{2}}$; this is possible since
$q_1$ is invertible modulo $q_2^2$, and $\gcd(c,q_2)=1$ because $cq_1\equiv p\pmod{q_2}$ and
$q_2\nmid p$. If $u\equiv c\pmod {q_2^{2}}$ then $q_2\mid p-uq_1$, so $v=(p-uq_1)/q_2$ is an
integer, and $p-uq_1\equiv -q_1q_2 \pmod {q_2^{2}}$ gives $v\equiv -q_1\not\equiv0\pmod{q_2}$.

Now put $R=q_1q_2^{2}$ and, using the Chinese remainder theorem, choose $\rho$ with $0\le\rho<R$,
\[
  \rho\equiv1 \pmod{q_1},\qquad \rho\equiv c \pmod{q_2^{2}} .
\]
Every $u\equiv\rho\pmod R$ then satisfies $q_1\nmid u$ automatically, as well as
$u\equiv c\pmod{q_2^2}$. Let $M=\lfloor p/q_1\rfloor$ and
\[
  S=\{\,u<M:\ u\equiv\rho \pmod R\,\},\qquad T=\#S .
\]
Since $S$ consists of the integers $\rho+Rt$ with $0\le t<T$, and since $M\ge p/q_1-1$ and
$R=q_1q_2^{2}$,
\begin{equation}\label{eq:count}
  T\ \ge\ \frac MR-1\ \ge\ \frac{p}{q_1^{2}q_2^{2}}-2 .
\end{equation}
For $1\le k\le g$ the sums $\frac1g\sum_{k=1}^{g}\psi^{k}(q_1u)$ detect the condition
$\psi(q_1u)=1$: the sum is $1$ when $\psi(q_1u)=1$ and $0$ otherwise, because $\psi(q_1u)$ is then
a $g$-th root of unity $\ne1$. Hence
\[
  \#\{u\in S:\ \psi(q_1u)=1\}\ =\ \frac1g\sum_{k=1}^{g}\ \sum_{u\in S}\psi^{k}(q_1u)
  \ =\ \frac{T}g+\frac1g\sum_{k=1}^{g-1}\ \sum_{u\in S}\psi^{k}(q_1u),
\]
the term $k=g$ contributing $T$ because $\psi^{g}=1$ and $\gcd(q_1u,p)=1$ for $u\in S$
(indeed $0<u<p$ and $q_1\ne p$). For $1\le k\le g-1$ the character $\psi^{k}$ is non-principal,
and the substitution $u=\rho+Rt$ turns the inner sum into a character sum over an interval:
writing $b=q_1R$, which is invertible modulo $p$, and $c_0=b^{-1}q_1\rho$ in $\Z/p\Z$, we have
$q_1(\rho+Rt)=b(c_0+t)$, so that
\[
  \Bigl|\sum_{u\in S}\psi^{k}(q_1u)\Bigr|
  =\Bigl|\psi^{k}(b)\sum_{t<T}\psi^{k}(c_0+t)\Bigr|
  \ \le\ p^{1/2}\bigl(1+\log p\bigr)
\]
by \eqref{eq:PV}. Combining the last two displays with \eqref{eq:count},
\[
  g\cdot\#\{u\in S:\ \psi(q_1u)=1\}\ \ge\ \frac{p}{q_1^{2}q_2^{2}}-2-(g-1)p^{1/2}(1+\log p)\ >\ 0
\]
by \eqref{eq:gcd-condition}. So there is $u\in S$ with $\psi(q_1u)=1$; then $u>0$ since $q_1\nmid
u$, $v=(p-uq_1)/q_2$ is a positive integer with $q_2\nmid v$, and $uq_1$ is an $n$-th power
residue modulo $p$. This is \eqref{eq:target}, and Lemma~\ref{lem:heilbronn-applies} shows that
$K_f$ is not norm-Euclidean.
\end{proof}

Combining Theorem~\ref{thm:gcd} with Proposition~\ref{prop:two-rootless} gives the asymptotic form
announced in the introduction. Because the character sums in the proof are taken modulo the prime
$p$ alone, their bound $p^{1/2}(1+\log p)$ does not grow with the auxiliary primes; only the main
term $p/(q_1^2q_2^2)$ shrinks. The admissible range of auxiliary primes is therefore a
\emph{power} of $p$, not a power of $\log p$.

\begin{corollary}\label{cor:gcd-density}
Let $n\ge2$, let $m\ge1$ be coprime to $n$, let $p$ be a prime and let $Y\ge2$ satisfy
\begin{equation}\label{eq:Y-condition}
  8\,n\,Y^{4}\ \le\ \lfloor p^{1/4}\rfloor .
\end{equation}
Then, with $\hat\varepsilon(p,Y)=(1+\pi(Y))(3/4)^{\pi(Y)}$,
\[
  \liminf_{X\to\infty}\delta_{p,m,n}(X)\ \ge\ 1-\hat\varepsilon(p,Y).
\]
One may take $Y=\lfloor p^{1/17}\rfloor$ as soon as $p\ge(8n)^{68}$, so that
$\hat\varepsilon(p,Y)\to0$ as $p\to\infty$ for each fixed $n$.
\end{corollary}

\begin{proof}
Let $Q$ be the set of primes $q\le Y$ and let $\mathcal F^{*}$ be the set of
$f\in\mathcal F_{p,m,n}$ having no root modulo $q$ for at least two $q\in Q$. By
Proposition~\ref{prop:two-rootless} the density of $\mathcal F^{*}$ inside $\mathcal F_{p,m,n}$ is
at least $1-\hat\varepsilon(p,Y)$, since $\#Q=\pi(Y)$. Fix $f\in\mathcal F^{*}$ and primes
$q_1<q_2\le Y$ modulo which $f$ has no root; both are different from $p$, because
$Y^{4}\le p^{1/4}$ forces $Y<p$. Writing $g=\gcd(p-1,n)$ we have $g\le n$ and
$q_1^{2}q_2^{2}\le Y^{4}$, so \eqref{eq:Y-condition} gives
$8\,g\,q_1^{2}q_2^{2}\le\lfloor p^{1/4}\rfloor$. Put $t=p^{1/4}$. Then $\sqrt p=t^{2}$ and, since
$\log p=4\log t\le 4(t-1)$,
\[
  (g-1)\,p^{1/2}(1+\log p)\ \le\ 4g\,t^{3}
  \qquad\text{while}\qquad
  \frac{p}{q_1^{2}q_2^{2}}=\frac{t^{4}}{q_1^{2}q_2^{2}}\ \ge\ 8g\,t^{3},
\]
the last step because $q_1^2q_2^2\le t/(8g)$. As $t\ge8$ and $g\ge1$ we get $4gt^{3}+2<8gt^{3}$,
which is \eqref{eq:gcd-condition}. Theorem~\ref{thm:gcd} therefore applies to every
$f\in\mathcal F^{*}$.

For the final assertion, $Y=\lfloor p^{1/17}\rfloor$ gives $Y^{4}\le p^{4/17}$, and
$8np^{4/17}\le p^{1/4}$ as soon as $8n\le p^{1/68}$, that is $p\ge(8n)^{68}$; and
$\pi(\lfloor p^{1/17}\rfloor)\to\infty$.
\end{proof}

\begin{remark}
In \cite{HMT} the divisibility $q_2\mid v$ is instead repaired afterwards, by replacing
$(u,v)$ with $(u+q_2,v-q_1)$ or with $(u+2q_2,v-2q_1)$. This does restore $q_2\nmid v$ and
$q_1\nmid u$, but it also changes $a=uq_1$ into $a=(u+q_2)q_1$ or $(u+2q_2)q_1$, and there is no
reason for the new value to remain an $n$-th power residue modulo $p$, which is what Heilbronn's
criterion needs. Detecting the condition $q_2\nmid v$ directly, as a congruence modulo $q_2^{2}$,
avoids the issue.
\end{remark}

\begin{remark}
Imposing the two congruence conditions on $u$ as the single congruence $u\equiv\rho\pmod{q_1q_2^2}$
is what keeps the estimate explicit. Detecting them by characters modulo $q_1$ and modulo $q_2^2$
instead leads to sums of characters modulo $M=pq_1q_2^{2}$, most of which are imprimitive; bounding
those requires passing to the inducing primitive character and costs a factor $2^{\omega(M)}$, and
the resulting condition on $g$ carries an unspecified absolute constant; worse, the bound then
grows with the auxiliary primes, which forces them to be of size $O((\log p)^{1/5})$. With a
single progression the only characters that occur are the $\psi^{k}$ modulo the prime $p$, for
which \eqref{eq:PV} applies directly and with a bound independent of $q_1,q_2$; the auxiliary
primes may then be as large as a fixed power of $p$, as in Corollary~\ref{cor:gcd-density}.
\end{remark}

\begin{corollary}
Fix $n\ge2$. A positive proportion of the monic integer polynomials of degree $n$ satisfying an
Eisenstein--Dumas condition fail to generate norm-Euclidean fields.
\end{corollary}

\begin{proof}
For odd $n$ this is Corollary~\ref{cor:positive}. For even $n$, choose a prime
$p\ge(8n)^{68}$ with $p\equiv-1\pmod{2n}$, which exists by Dirichlet's theorem; then
$\gcd(p-1,n)=2$, and Corollary~\ref{cor:gcd-density} with $Y=\lfloor p^{1/17}\rfloor$ applies and
gives a positive lower density inside $\mathcal F_{p,m,n}$, a family of positive density by
Theorem~\ref{thm:main-density-family}.
\end{proof}

\begin{thebibliography}{99}

\bibitem{Dubickas} A.~Dubickas, \emph{Polynomials irreducible by Eisenstein's criterion},
Appl. Algebra Engrg. Comm. Comput. \textbf{14} (2003), 127--132.

\bibitem{Dumas} G.~Dumas, \emph{Sur quelques cas d'irr\'eductibilit\'e des polyn\^omes \`a
coefficients rationnels}, J. Math. Pures Appl. \textbf{2} (1906), 191--258.

\bibitem{Egami} S.~Egami, \emph{On finiteness of the numbers of Euclidean fields in some classes
of number fields}, Tokyo J. Math. \textbf{7} (1984), 183--198.

\bibitem{EK} P.~Erd\H{o}s and C.~Ko, \emph{Note on the Euclidean algorithm},
J. London Math. Soc. \textbf{13} (1938), 3--8.

\bibitem{H1} H.~Heilbronn, \emph{On Euclid's algorithm in cubic self-conjugate fields},
Proc. Cambridge Philos. Soc. \textbf{46} (1950), 377--382.

\bibitem{H2} H.~Heilbronn, \emph{On Euclid's algorithm in cyclic fields},
Canad. J. Math. \textbf{3} (1951), 257--268.

\bibitem{HMT} A.~Hibbler, K.~J.~McGown and E.~Trevi\~no, \emph{Polynomial densities and
Heilbronn's criterion}, arXiv:2512.16220.

\bibitem{Lemmermeyer} F.~Lemmermeyer, \emph{The Euclidean algorithm in algebraic number fields},
Exposition. Math. \textbf{13} (1995), 385--416.

\bibitem{Polya} G.~P\'olya, \emph{\"Uber die Verteilung der quadratischen Reste und Nichtreste},
Nachr. Akad. Wiss. G\"ottingen (1918), 21--29.

\bibitem{Vinogradov} I.~M.~Vinogradov, \emph{Sur la distribution des r\'esidus et des non-r\'esidus
des puissances}, J. Phys.-Math. Soc. Perm \textbf{1} (1918), 94--98.

\end{thebibliography}

\end{document}
